\clearpage
\begin{examplebox}
\textbf{例1（顶点切线：十秒口算）。}

$H:x^2/9-y^2/4=1$，在右顶点作切线，切线与渐近线交于 $A,B$，求面积。

\textbf{解：} $a=3,b=2$；顶点切线为 $x=3$，交点为
$A=(3,2)$、$B=(3,-2)$，故
$S_{\triangle AOB}=\frac12\times4\times3=\boxed{6}$。

\medskip
\textbf{例2（反向应用：由面积和长度乘积求半轴）。}

某双曲线 $a>b>0$，其切线与两渐近线交于 $A,B$，已知
$S_{\triangle AOB}=6$，$OA\cdot OB=13$，求 $a,b$。

\textbf{解：} $ab=6$，$a^2+b^2=13$。
因此 $(a+b)^2=25$、$(a-b)^2=1$，又 $a>b>0$，
故 $\boxed{a=3,\ b=2}$。

\end{examplebox}
\clearpage
\begin{examplebox}
\textbf{例3（非顶点：交点数值检验中点和面积）。}

同一双曲线，过 $P=(5,8/3)$ 作切线，求 $A,B$ 及 $S_{\triangle AOB}$。
\begin{center}
% Exact example: a=3,b=2,r=3; OAP and OPB have equal area 3.
\begin{tikzpicture}[x=.87cm,y=.87cm,>=stealth,every node/.style={font=\small}]
\coordinate (O) at (0,0);\coordinate (A) at (9,6);
\coordinate (B) at (1,-.666667);\coordinate (P) at (5,2.666667);
\fill[orange!16] (O)--(A)--(P)--cycle;
\fill[red!10] (O)--(P)--(B)--cycle;
\draw[gray!65,->] (-.5,0)--(10.2,0) node[right] {$x$};
\draw[gray!65,->] (0,-1.3)--(0,6.8) node[above] {$y$};
\draw[blue!65!black,densely dashed] (-.4,-.266667)--(9.6,6.4);
\draw[blue!65!black,densely dashed] (-.4,.266667)--(1.8,-1.2);
\draw[teal!70!black,thick,domain=-.5:1.82,samples=90,variable=\t]
 plot ({1.5*(exp(\t)+exp(-\t))},{exp(\t)-exp(-\t)});
\draw[red!75!black,very thick] ($(B)!-.04!(A)$)--($(A)!-.04!(B)$);
\draw[orange!90!black,thick] (A)--(O)--(B);
\draw[black,thick] (O)--(P);
\fill (O) circle(1.5pt) node[below left=5pt] {$O$};
\fill (A) circle(1.5pt) node[above left] {$A(9,6)$};
\fill (B) circle(1.5pt) node[below=10pt] {$B(1,-\frac23)$};
\fill (P) circle(1.5pt) node[right=6pt] {$P(5,\frac83)$};
\node at (7.4,1.3) {$S_{OAP}=3$};
\node at (7.4,.3) {$S_{OPB}=3$};
\end{tikzpicture}

\end{center}

\textbf{解：} $P\in H$，且 $u=5/3,\ v=4/3$。代入交点速式
\[
A=\left(\frac3{u-v},\frac2{u-v}\right)=(9,6),\qquad
B=\left(\frac3{u+v},-\frac2{u+v}\right)
=\left(1,-\frac23\right).
\]
验算中点 $(A+B)/2=(5,8/3)=P$；且
\[
S=\frac12\left|9\left(-\frac23\right)-1\cdot6\right|=\boxed6.
\]
两条渐近线上的交点位置变化很大，但面积没有变化。

\medskip
\textbf{变式（迁移）：}仍在 $x^2/9-y^2/4=1$ 上取任意点切线。
若该切线与两渐近线交点是 $A,B$，求
$S_{\triangle OAP}$、$S_{\triangle OBP}$。

\textbf{答案：}因为 $P$ 是 $AB$ 的中点，$OP$ 为中线，
\[
\boxed{S_{\triangle OAP}=S_{\triangle OBP}=\frac{ab}{2}=3}.
\]
\end{examplebox}
